\subsection{Konsequenzen: Vom Kathodenstrahl zum Quantenteilchen}
\label{subsec:konsequenzen-kathodenstrahl-thomson-millikan-photoeffekt}

Die in diesem Kapitel behandelten Experimente markieren einen tiefen
Umbruch in der Physik. Aus zunächst rätselhaften Leuchterscheinungen in
Entladungsröhren entwickelte sich Schritt für Schritt ein neues Bild der
Materie: Elektrizität war nicht nur ein kontinuierliches Phänomen, sondern
an kleine, negativ geladene Teilchen gebunden.

Die Kathodenstrahlen zeigten zunächst, dass in evakuierten Röhren Strahlen
auftreten, die von elektrischen und magnetischen Feldern abgelenkt werden
können. Damit wurde klar, dass es sich nicht um gewöhnliches Licht handelt,
sondern um etwas, das Ladung und Trägheit besitzt.

Joseph John Thomson\index{Thomson, Joseph John} ging den entscheidenden
Schritt weiter. Durch die Ablenkung der Kathodenstrahlen in elektrischen und
magnetischen Feldern konnte er das Verhältnis von Ladung zu Masse bestimmen:

\[
\frac{e}{m}.
\]

Damit zeigte er, dass die Träger der Kathodenstrahlen wesentlich leichter
sein mussten als Atome. Das Elektron\index{Elektron} erschien erstmals als
eigenständiges subatomares Teilchen. Die Vorstellung vom unteilbaren Atom
war damit endgültig erschüttert.

Robert Andrews Millikan\index{Millikan, Robert Andrews} ergänzte dieses Bild
durch sein Öltröpfchenexperiment\index{Öltröpfchenexperiment}. Er bestimmte
nicht nur die Größe der Elementarladung\index{Elementarladung}, sondern
zeigte auch, dass elektrische Ladung nur in ganzzahligen Vielfachen einer
kleinsten Einheit auftritt:

\[
q = n e,
\qquad n \in \mathbb{Z}.
\]

Damit wurde deutlich: Ladung ist quantisiert. Elektrizität besitzt also eine
körnige Struktur. Aus dem Verhältnis \(e/m\) von Thomson und der Elementarladung
\(e\) von Millikan ließ sich schließlich auch die Masse des Elektrons bestimmen:

\[
m_e = \frac{e}{e/m}.
\]

Das Elektron war nun nicht mehr nur ein hypothetischer Träger elektrischer
Erscheinungen, sondern ein messbares Teilchen mit Ladung und Masse.

Eine weitere entscheidende Konsequenz ergab sich aus dem photoelektrischen
Effekt\index{Photoelektrischer Effekt}. Dieser zeigte, dass Elektronen durch
Licht aus Metalloberflächen herausgelöst werden können, jedoch nicht in der
Weise, wie es die klassische Wellentheorie erwarten ließ. Entscheidend war
nicht allein die Intensität des Lichts, sondern vor allem seine Frequenz.

Albert Einstein\index{Einstein, Albert} deutete dieses Verhalten durch die
Annahme, dass Licht seine Energie in einzelnen Quanten überträgt. Die Energie
eines solchen Lichtquants ist

\[
E = h \nu.
\]

Damit entstand eine neue Verbindung zwischen Licht und Elektron: Licht konnte
nicht nur als elektromagnetische Welle verstanden werden, sondern musste in
bestimmten Experimenten auch als quantisierte Energieübertragung beschrieben
werden. Das Elektron wurde dadurch zu einem zentralen Prüfstein der entstehenden
Quantenphysik.

\begin{DidacticBox}[Vom klassischen Teilchen zum Quantenteilchen]
	\label{box:vom-klassischen-teilchen-zum-quantenteilchen}
	Die Experimente dieses Kapitels zeigen drei entscheidende Schritte:
	\begin{itemize}
		\item Kathodenstrahlen machten sichtbar, dass bewegte Ladungsträger existieren.
		\item Thomson bestimmte das Verhältnis von Ladung zu Masse und identifizierte
		das Elektron als subatomares Teilchen.
		\item Millikan bestimmte die Elementarladung und zeigte die Quantisierung der
		elektrischen Ladung.
		\item Der photoelektrische Effekt verband das Elektron mit der Quantennatur
		des Lichts.
	\end{itemize}
	Damit wurde das Elektron vom klassischen Ladungsträger zu einem Schlüsselobjekt
	der modernen Physik.
\end{DidacticBox}

Diese Entwicklung hatte weitreichende Folgen. Das Atom konnte nicht länger als
unstrukturierte, unteilbare Einheit betrachtet werden. Es musste einen inneren
Aufbau besitzen. Zugleich wurde klar, dass elektrische Ladung, Lichtabsorption
und Energieübertragung nicht beliebig kontinuierlich verlaufen, sondern durch
diskrete Einheiten bestimmt sein können.

Damit führt dieses Kapitel direkt zum nächsten großen Schritt: Wenn Elektronen
Bestandteile der Materie sind, dann muss auch das Atommodell neu gedacht werden.
Die Frage lautet nun nicht mehr nur, ob Elektronen existieren, sondern wie sie
im Atom angeordnet sind und welche Rolle sie für die Struktur der Materie
spielen.